\documentclass[10pt,a4paper]{amsart}
\usepackage{amsmath,amssymb,amsthm,graphicx}
% Retained minimal preamble for the Stacks Project chapter topology.tex (Topology).
% The Stacks Project uses one single shared preamble for every chapter; the snapshot stages it
% identically as sources/stacks-topology/preamble.tex, sources/stacks-chow/preamble.tex and
% sources/stacks-etale-cohomology/preamble.tex.  The block below is copied mechanically, line for
% line, from that shared preamble: the Xy-pic setup (lines 14-18) and the native theorem
% declarations together with the mathematical macros (lines 154-217), plus the project's own
% reference/slogan/history display environments (lines 8-12, via the verbatim package).
% Nothing was added, removed or redefined; the wrapper class and the amsmath/amssymb/amsthm
% base packages are supplied by the assembler.
% Chapter: topology.tex -- Topology.
\usepackage[all]{xy}

% We use 2cell for 2-commutative diagrams.
\xyoption{2cell}

% Theorem environments.
%
\theoremstyle{plain}
\newtheorem{theorem}[subsection]{Theorem}
\newtheorem{proposition}[subsection]{Proposition}
\newtheorem{lemma}[subsection]{Lemma}

\theoremstyle{definition}
\newtheorem{definition}[subsection]{Definition}
\newtheorem{example}[subsection]{Example}
\newtheorem{exercise}[subsection]{Exercise}
\newtheorem{situation}[subsection]{Situation}

\theoremstyle{remark}
\newtheorem{remark}[subsection]{Remark}
\newtheorem{remarks}[subsection]{Remarks}

\numberwithin{equation}{subsection}

% Macros
%
\def\lim{\mathop{\mathrm{lim}}\nolimits}
\def\colim{\mathop{\mathrm{colim}}\nolimits}
\def\Spec{\mathop{\mathrm{Spec}}}
\def\Hom{\mathop{\mathrm{Hom}}\nolimits}
\def\Ext{\mathop{\mathrm{Ext}}\nolimits}
\def\SheafHom{\mathop{\mathcal{H}\!\mathit{om}}\nolimits}
\def\SheafExt{\mathop{\mathcal{E}\!\mathit{xt}}\nolimits}
\def\Sch{\mathit{Sch}}
\def\Mor{\mathop{\mathrm{Mor}}\nolimits}
\def\Ob{\mathop{\mathrm{Ob}}\nolimits}
\def\Sh{\mathop{\mathit{Sh}}\nolimits}
\def\NL{\mathop{N\!L}\nolimits}
\def\CH{\mathop{\mathrm{CH}}\nolimits}
\def\proetale{{pro\text{-}\acute{e}tale}}
\def\etale{{\acute{e}tale}}
\def\QCoh{\mathit{QCoh}}
\def\Ker{\mathop{\mathrm{Ker}}}
\def\Im{\mathop{\mathrm{Im}}}
\def\Coker{\mathop{\mathrm{Coker}}}
\def\Coim{\mathop{\mathrm{Coim}}}

% Boxtimes
%
\DeclareMathSymbol{\boxtimes}{\mathbin}{AMSa}{"02}

%
% Macros for moduli stacks/spaces
%
\def\QCohstack{\mathcal{QC}\!\mathit{oh}}
\def\Cohstack{\mathcal{C}\!\mathit{oh}}
\def\Spacesstack{\mathcal{S}\!\mathit{paces}}
\def\Quotfunctor{\mathrm{Quot}}
\def\Hilbfunctor{\mathrm{Hilb}}
\def\Curvesstack{\mathcal{C}\!\mathit{urves}}
\def\Polarizedstack{\mathcal{P}\!\mathit{olarized}}
\def\Complexesstack{\mathcal{C}\!\mathit{omplexes}}
% \Pic is the operator that assigns to X its picard group, usage \Pic(X)
% \Picardstack_{X/B} denotes the Picard stack of X over B
% \Picardfunctor_{X/B} denotes the Picard functor of X over B
\def\Pic{\mathop{\mathrm{Pic}}\nolimits}
\def\Picardstack{\mathcal{P}\!\mathit{ic}}
\def\Picardfunctor{\mathrm{Pic}}
\def\Deformationcategory{\mathcal{D}\!\mathit{ef}}

% --- appended 2026-10-12 (Chief): the project own definitions of the reference/slogan/history
% display environments, copied mechanically from the shared Stacks preamble lines 8-12
% (\usepackage{verbatim} plus the three \newenvironment lines, which the chapter preamble keeps
% verbatim from the project preamble). Nothing added, removed or redefined: these environments
% carry no mathematical content and the project itself typesets them as comments. Without these
% three lines a counted statement whose own text uses them fails with Environment slogan undefined.
\usepackage{verbatim}
\newenvironment{reference}{\comment}{\endcomment}
\newenvironment{slogan}{\comment}{\endcomment}
\newenvironment{history}{\comment}{\endcomment}
\title{A quasi-compact locally closed subset with retrocompact complement in a nice basis is constructible}
\author{The Stacks Project Authors}
\date{stacks-project:topology.tex:local-snapshot-20261012; CC BY 4.0}
\begin{document}
\maketitle
\noindent\textit{Editorial note.} All delivered mathematical text is the author's own, copied from the pinned Stacks Project chapter source sources/stacks-topology/pinned-original-source.tex; no statement or proof step was rewritten, reordered or removed.\par\medskip
\noindent\textit{Editorial note.} Preamble provenance: the Stacks Project keeps one single shared preamble for every chapter, staged here as sources/stacks-topology/preamble.tex; this card uses retained-preamble.tex, which copies that shared preamble mechanically, including its reference/slogan/history display environments. Nothing was added, removed or redefined; the wrapper class and the base packages are supplied by the assembler.\par\medskip
\noindent\textit{Editorial note.} Bibliography provenance: the chapter carries no chapter-local bibliography and neither the counted statement nor its proof contains a citation, so the card carries no bibliography and no entry was reconstructed.\par\medskip
\noindent\textit{Editorial note.} Selection provenance: the card was taken from the staged candidate list for this chapter, candidate lemma-locally-closed-constructible-image; the statement and proof spans are that candidate's own statement\_lines and proof\_lines, checked against the pinned chapter so that each span opens and closes a real, uncommented environment.\par\medskip
\noindent\textit{Allowed standard prerequisites (editorial).}\begin{itemize}
\item The notions of a retrocompact subset and of a constructible subset, and the hypothesis that the space has a basis of quasi-compact opens closed under finite intersections. Set up in the author's section on constructible sets, immediately before the counted statement. Fixes the hypotheses on the space and on the subset, and the conclusion.
\end{itemize}\medskip
\begin{lemma}
\label{lemma-locally-closed-constructible-image}
Let $X$ be a quasi-compact topological space having a basis consisting of
quasi-compact opens such that the intersection of any two
quasi-compact opens is quasi-compact.
Let $T \subset X$ be a locally closed subset
such that $T$ is quasi-compact and $T^c$ is retrocompact in $X$.
Then $T$ is constructible in $X$.
\end{lemma}

\begin{proof}
Note that $T$ is quasi-compact and open in $\overline{T}$.
Using our basis of quasi-compact opens we can write
$T = U \cap \overline{T}$ where $U$ is quasi-compact open in $X$.
Then $V = U \setminus T = U \cap T^c$ is retrocompact in $U$ as $T^c$
is retrocompact in $X$. Hence $V$ is quasi-compact. Since
the intersection of any two quasi-compact opens is quasi-compact
any quasi-compact open of $X$ is retrocompact. Thus $T = U \cap V^c$
with $U$ and $V = U \setminus T$ retrocompact opens of $X$.
A fortiori, $T$ is constructible in $X$.
\end{proof}

\end{document}
